\documentclass[11pt,a4paper]{article}

\usepackage[utf8]{inputenc}
\usepackage[T1]{fontenc}
\usepackage[margin=2.4cm]{geometry}
\usepackage{booktabs}
\usepackage{longtable}
\usepackage{array}
\usepackage{xcolor}
\usepackage{graphicx}
\usepackage{listings}
\usepackage{enumitem}
\usepackage[hidelinks]{hyperref}

\definecolor{codebg}{gray}{0.96}
\definecolor{rule}{gray}{0.80}

\lstset{
  basicstyle=\ttfamily\scriptsize,
  backgroundcolor=\color{codebg},
  breaklines=true,
  breakatwhitespace=false,
  columns=fullflexible,
  frame=single,
  rulecolor=\color{rule},
  xleftmargin=0pt,
  framexleftmargin=4pt,
  framexrightmargin=4pt,
  aboveskip=8pt,
  belowskip=8pt,
  literate={->}{{$\rightarrow$}}2
}

\newcommand{\id}{24B0073}
\setlength{\parskip}{5pt}
\setlength{\parindent}{0pt}

\title{\textbf{Food Web Council}\\[4pt]
\large A multi-agent learning environment for food-chain disruption}
\author{Roll number \id{} \quad\textbar\quad ET 617 End-Semester Examination}
\date{}

\begin{document}
\maketitle

%====================================================================
\section*{Identity card}
\addcontentsline{toc}{section}{Identity card}

\begin{tabular}{@{}p{3.6cm}p{12cm}@{}}
\toprule
\textbf{Field} & \textbf{Value}\\
\midrule
P --- Problem      & \textbf{P3} Multi-agent collaborative learning environment\\
T --- Topic        & \textbf{P3-T4} Food chains: effect of removing a species (Cl.\ 7)\\
L --- Learner      & \textbf{L5} Learner repeating the topic after failing it once; high anxiety, low self-efficacy\\
C --- Constraint   & \textbf{C1} Low bandwidth\\
ID                 & \textbf{\id}\\
\midrule
Live application   & \url{https://et617-24b0073.vercel.app}\\
Repository         & \url{https://github.com/Rajas129/et617-endsem-24b0073}\\
\bottomrule
\end{tabular}

\vspace{6pt}
\textit{Note on the topic code.} In the P3 column of Section 4.1, ``Food chains: effect of
removing a species'' is \textbf{T4}; T3 in that column is series versus parallel circuits.
The card is therefore declared as P3-T4.

%====================================================================
\section{D1 \quad Educational problem}

Class 7 learners can read a single food chain from left to right, but they treat it as a
\emph{line} rather than a \emph{web}. Asked what happens when one species is removed, they
predict an effect only on the animal immediately adjacent to it, and they assume that removing
a predator is straightforwardly good for its prey. Both ideas survive ordinary teaching because
a textbook diagram is static: there is no way to pull a species out and watch what happens three
arrows away. A learner who has already failed this unit once carries a second problem on top of
the first --- she has learned to expect being told she is wrong, so any design that feels like a
test stops her before the content does.

\subsection*{Learning outcomes}

\begin{enumerate}[leftmargin=*,itemsep=2pt]
  \item \textbf{Predict} the effect of removing one species on at least two species that are
        \emph{not} adjacent to it in the web.
  \item \textbf{Explain} why a predator may survive the loss of one prey species, with reference
        to its alternate feeding links.
  \item \textbf{Compare} two conflicting predictions and \textbf{identify} the assumption that
        separates them.
\end{enumerate}

All three are observable in the application's own logs: outcome 1 from the species she ticks,
outcome 2 from her free-text explanation, outcome 3 from whether she edits the deciding arrow.

\subsection*{Context of use}

On her own, away from the classroom in which she failed, in a web browser on a laptop or
shared computer over an unreliable connection. She may repeat the activity as often as she
likes before the re-test.

\textbf{Constraint scope.} The card's constraint is \textbf{C1 (low bandwidth)}. Session length
is C4 and one-handed mobile layout is C3; neither is this card's constraint, and nothing in this
design is justified by them. The application has no timer and makes no claim about how long a
sitting should be.

%====================================================================
\section{D2 \quad Learner persona}

\textbf{Seed: L5} --- a learner repeating the topic after failing it once; high anxiety, low
self-efficacy. The persona below is expanded from that seed alone.

\textbf{Priya Kamble, 13.} Repeating Class 7 science after failing the ecosystems unit last term.

\begin{longtable}{@{}p{3.3cm}p{12.3cm}@{}}
\toprule
\textbf{Attribute} & \textbf{Detail}\\
\midrule
\endhead
Prior knowledge &
Not a beginner, and this matters. She has already been taught food chains once. She names
producer, consumer and decomposer correctly and reads a single chain left to right. What did not
survive the first attempt is the \emph{web} --- two arrows arriving at one animal --- which is
precisely what she was tested on and failed.\\
\addlinespace
Misconception 1 &
\textbf{Linear chain.} Removing a species affects only the animal directly above it.
\textit{``Frogs gone, so snakes. Grass is far away, grass is fine.''}\\
\addlinespace
Misconception 2 &
\textbf{Predator removal is good.} Removing a predator simply means more prey, and more is good.
Misses the overgrazing crash that follows the boom.\\
\addlinespace
Language &
Simple English. Under stress she re-reads the same sentence and stops, so every sentence she is
shown is short and plainly worded.\\
\addlinespace
Device &
A web browser on a laptop or shared computer, on a connection she cannot rely on.\\
\addlinespace
Time &
Unbounded and self-paced. Nothing in the application is timed or capped.\\
\addlinespace
Motivation &
Low and defensive. She is not trying to learn ecology; she is trying not to fail again. The
activity has to be worth doing before it is worth finishing.\\
\addlinespace
Affect &
The defining trait. She has already been publicly wrong about this topic. She freezes at scores
and red marks, will not type into a box she believes is graded, and reads a confident correction
as evidence that she is bad at science. Being wrong a \emph{second} time is what she is
protecting herself against, so the design must make being wrong survivable before it can make it
useful.\\
\bottomrule
\end{longtable}

%====================================================================
\section{D3 \quad Educational content}

\subsection*{The concept explanation}

The learner first meets three short ideas, before any activity: that a food chain shows who eats
whom; that the arrow points \emph{from the food to whoever eats it}; and that real life is not a
single line, because most animals eat more than one thing. The third idea is the one the whole
activity turns on and it is stated plainly up front: ``that makes a \emph{web}, not a chain ---
and it is why taking one animal out can change something far away from it.''

She then names any animal she likes and a food web is built around it. The worked example below
uses the canonical web, whose figures are the verified output of the simulator.

\subsection*{The web}

Seven species. The snake eats frogs (0.55 of its diet) \emph{and rats} (0.45). That one split is
the pedagogical centre of the activity: it is what allows two agents to disagree honestly.

\begin{center}
\begin{tabular}{@{}llr@{}}
\toprule
\textbf{Species} & \textbf{Role} & \textbf{Baseline}\\
\midrule
Grass & producer & 100\\
Grasshopper & consumer (eats grass) & 60\\
Rat & consumer (eats grass) & 40\\
Frog & consumer (eats grasshopper) & 50\\
Snake & consumer (eats frog 0.55, rat 0.45) & 40\\
Hawk & consumer (eats snake 0.70, rat 0.30) & 20\\
Decomposers & decomposer & 30\\
\bottomrule
\end{tabular}
\end{center}

\subsection*{Verified results}

These are simulator output, not estimates. With nothing removed, every species holds exactly at
its baseline for all 12 generations, so any change the learner sees is attributable to something
she did.

\begin{center}
\begin{tabular}{@{}lrrll@{}}
\toprule
\textbf{Species} & \textbf{Start} & \textbf{End} & \textbf{Direction} & \textbf{Hops from frog}\\
\midrule
\multicolumn{5}{@{}l}{\textit{Remove the frog}}\\
Grasshopper & 60 & 188.4 & boomed & 1\\
Snake & 40 & 18.6 & fell & 1\\
\textbf{Grass} & \textbf{100} & \textbf{57.9} & \textbf{fell} & \textbf{2}\\
\textbf{Hawk} & \textbf{20} & \textbf{15.5} & \textbf{fell} & \textbf{2}\\
Rat & 40 & 36.6 & steady & 2\\
\midrule
\multicolumn{5}{@{}l}{\textit{Remove the frog, with the Rat $\rightarrow$ Snake arrow deleted}}\\
Snake & 40 & \textbf{0} & \textbf{died out} & 1\\
\midrule
\multicolumn{5}{@{}l}{\textit{Remove the snake} (targets misconception 2)}\\
Rat & 40 & 91.2 & boomed & 1\\
Frog & 50 & 70.3 & rose & 1\\
\textbf{Grasshopper} & \textbf{60} & \textbf{11.1} & \textbf{fell} & \textbf{2}\\
Hawk & 20 & 12.9 & fell & 1\\
\bottomrule
\end{tabular}
\end{center}

The bold rows are the teaching. Grass and hawk are two hops from the frog and still move, which
contradicts misconception 1. Removing the snake booms the rats \emph{and crashes the
grasshoppers}, which contradicts misconception 2 using the learner's own experiment rather than
a correction.

\subsection*{Five tasks in four formats}

\begin{longtable}{@{}p{0.5cm}p{2.9cm}p{8.4cm}p{3.4cm}@{}}
\toprule
\textbf{\#} & \textbf{Format} & \textbf{Task} & \textbf{Targets}\\
\midrule
\endhead
1 & Multi-select & ``Tick everyone who changes.'' Correct: Grasshopper, Snake, \emph{Grass},
    \emph{Hawk} & Misconception 1\\
2 & Prediction commit & ``Do the snakes die out, drop but survive, hold steady, or rise?''
    Committed \emph{before} anything runs & Outcome 1\\
3 & Free text & ``The snakes survived. Why?'' Evaluated against the simulation & Outcome 2\\
4 & Model edit & ``Delete one arrow so that the snakes \emph{do} die out.'' & Outcome 3\\
5 & Adjudication & ``Mira and Aarav disagree. Who is right, and what would settle it?'' & Outcome 3\\
\bottomrule
\end{longtable}

\subsection*{Model answers}

\begin{description}[leftmargin=1.2cm,style=nextline,itemsep=2pt]
  \item[Task 1] Grasshopper, Snake, Grass and Hawk. Rat and Decomposers barely move.
  \item[Task 2] ``They will drop but survive.'' Snakes fall from 40 to 18.6 and hold.
  \item[Task 3] ``Because snakes eat rats as well as frogs, so they still had food.''
  \item[Task 4] Delete \texttt{Rat $\rightarrow$ Snake}. Snakes then fall from 40 to 0.
  \item[Task 5] ``Neither is right on its own --- it depends whether the snake has another food.
        Deleting that arrow and running it again would settle it.''
\end{description}

\subsection*{Feedback (written in code, so it cannot drift between runs)}

\textbf{For a correct response to task 1:}
\begin{quote}\itshape
``You got Grass and Hawk too. That's the hard one --- most people stop at the animal right next
to them. They changed because the whole web is joined up.''
\end{quote}

\textbf{When misconception 1 fires} (she ticks only the snake):
\begin{quote}\itshape
``Snake is right --- good. But let's keep going down the arrows. If the frogs are gone, who stops
being eaten? And then what happens to what they eat? Look at Grass and Hawk.''
\end{quote}
She is shown the cascade and never told the answer. The far-away change is then stated plainly:
\textit{``Grass and Hawk are not even touching the frog in the web, and they still changed.''}

\textbf{When misconception 2 fires} (she expects a predator's removal to simply help its prey):
\begin{quote}\itshape
``Grasshopper did boom --- up to 196 around step 8. But look at Grass: 100 $\rightarrow$ 58. More
grasshoppers was not simply good news. Why do you think Grass dropped?''
\end{quote}
The overshoot is the teacher. In the snake-removal run the same mechanism is starker: rats boom
to 91.2 while grasshoppers crash to 11.1.

%====================================================================
\section{D4 \quad Working application}

\textbf{URL:} \url{https://et617-24b0073.vercel.app} \qquad
\textbf{Identity token \id{}} appears in the footer of every screen and on the About screen.

\subsection*{Screen map (eleven connected states)}

\begin{longtable}{@{}p{0.5cm}p{2.6cm}p{5.2cm}p{6.6cm}@{}}
\toprule
\textbf{\#} & \textbf{Screen} & \textbf{Learner input} & \textbf{System output / transition}\\
\midrule
\endhead
1 & Sign in & username, password & session cookie; progress stored per user\\
2 & What is a food chain? & reads, continues & definition, arrow convention, worked 3-node example\\
3 & Pick an animal & types any animal & $\rightarrow$ building\\
4 & Building & waits & AI proposes structure; code computes all coefficients\\
5 & The web & clicks a species & what it eats / what eats it; ``take them out'' $\rightarrow$ predict\\
6 & Predict & multi-select, then a committed prediction + free text & $\rightarrow$ council\\
7 & Council & reads two agents & if they conflict, a Two Views card appears\\
8 & Two Views & chooses a side (or ``both wrong'') & $\rightarrow$ run the simulation\\
9 & Results & reads, explains in own words & cascade, sparklines, table, Dr.\ Sen's verdict\\
10 & Ask Mira & types any doubt & answer; may propose a web edit to try\\
11 & What we settled & reviews & accumulated findings, each with ``How we know''\\
\bottomrule
\end{longtable}

\subsection*{C1 (low bandwidth) --- how it is visibly satisfied}

\begin{longtable}{@{}p{5.2cm}p{10.4cm}@{}}
\toprule
\textbf{Requirement} & \textbf{Implementation}\\
\midrule
\endhead
Core activity works on a slow connection &
The population simulator is ordinary arithmetic running in the browser. Building the web,
removing a species, committing a prediction and running 12 generations all work with
\emph{no network at all}. The AI is commentary on top, never load-bearing.\\
\addlinespace
No media over 1\,MB &
Zero images, zero audio, zero video, zero web fonts. The food web and every chart are inline
SVG; the system font stack is used throughout. Measured first load: \textbf{120\,kB}.\\
\addlinespace
A clear fallback when an AI call fails &
A labelled banner plus a recorded conversation, while the simulator stays live. Three distinct
causes are distinguished rather than collapsed into ``offline'' (see D6).\\
\addlinespace
Visible to the examiner &
The footer carries a live byte counter (``41 KB used'') that ticks once per model round-trip,
beside the identity token.\\
\bottomrule
\end{longtable}

Further C1-driven decisions: at most two model calls per round rather than four; per-agent token
caps that also keep the reading level low; and the session persisted to the browser so a dropped
connection never loses her work.

\subsection*{Screenshots}

\textit{Insert captured screenshots of the live application here.}

% \begin{figure}[h]\centering
%   \includegraphics[width=0.48\textwidth]{screens/web.png}
%   \includegraphics[width=0.48\textwidth]{screens/twoviews.png}
%   \caption{The generated web, and the Two Views disagreement card.}
% \end{figure}

%====================================================================
\section{D5 \quad Two contrasting learner scenarios}

Both scenarios were run on the application. They differ at the first task and the system
responds differently from that point onward.

\subsection*{Scenario A --- the learner who responds correctly}

She ticks Grasshopper, Snake, Grass and Hawk, and predicts that the snakes will drop but survive.

\begin{enumerate}[leftmargin=*,itemsep=1pt]
  \item \textbf{Task 1 feedback.} ``You got Grass and Hawk too. That's the hard one --- most
        people stop at the animal right next to them.''
  \item \textbf{Council.} Aarav still disagrees with Mira (his stance is scheduled, not a
        reaction to her being right), so the Two Views card still appears. She now has to defend
        a position she already holds, which is a different cognitive task from forming one.
  \item \textbf{Simulation.} Snakes 40 $\rightarrow$ 18.6. Her prediction is confirmed.
  \item \textbf{Dr.\ Sen.} Confirms, then names the arrow the result depends on and invites the
        deciding experiment. \emph{Being right does not end the activity} --- she is still sent
        to test the assumption.
  \item \textbf{Settled.} ``A predator survives losing one prey only if it has another.''
\end{enumerate}

\subsection*{Scenario B --- the learner who shows misconception 1}

She ticks \emph{only} the snake --- the frog's immediate neighbour --- and predicts the snakes
will die out completely.

\begin{enumerate}[leftmargin=*,itemsep=1pt]
  \item \textbf{Task 1 feedback diverges immediately.} ``Snake is right --- good. But let's keep
        going down the arrows \ldots Look at Grass and Hawk.'' She is given something to watch,
        not a correction.
  \item \textbf{Council.} She chooses ``Aarav is right'', siding with her own misconception.
  \item \textbf{Simulation.} Snakes 40 $\rightarrow$ 18.6: her prediction does not hold. The
        indirect note fires --- \textit{``Grass and Hawk are not even touching the frog in the
        web, and they still changed.''} --- which Scenario A never sees, because she had already
        ticked them.
  \item \textbf{Dr.\ Sen.} ``So the snakes did not die out. They dropped to 19 and held on. The
        reason is one arrow: Rat $\rightarrow$ Snake. Delete that arrow and run it again.''
  \item \textbf{The resolution.} She deletes the arrow. Snakes fall to \textbf{0}.
        \begin{quote}\itshape
        ``With the Rat $\rightarrow$ Snake arrow the snakes ended at 19. Without it they ended at
        0. So Aarav was right after all --- but only in a web where the snakes have nothing else
        to eat. Neither of them was wrong about nature. They were picturing different webs. That
        one arrow was the whole argument, and you settled it yourself.''
        \end{quote}
  \item \textbf{Misconception state.} Her linear-chain confidence rises on the mis-tick and falls
        again once she identifies the deciding arrow, which raises the scaffolding level in
        between.
\end{enumerate}

\textbf{How the system responds differently:} different task-1 feedback; an indirect-effects note
that only the misconception path receives; a different ``who was right'' verdict; and an extra
required experiment before anything is recorded as settled.

\textit{Insert screenshots of both paths here.}

%====================================================================
\section{D6 \quad AI functionality specification}

Four AI features. For each: what it does, what it receives, what it emits, how the output is
checked before the learner sees it, and what happens when it fails.

\subsection*{Where the AI is --- and where it deliberately is not}

The population simulation contains \textbf{no AI}. It is arithmetic in the browser and it is the
ground truth against which AI claims are checked. This is the single most important design
decision in the system: it is what makes verification possible at all.

\subsection*{Feature 1 --- Building the food web}

\textbf{Does:} turns an animal the learner names into a 3--5 stage food web.

\textbf{Receives:} the animal's name, inside this prompt (abridged):

\begin{lstlisting}
Build a simple food web around "<animal>" for a Class 7 science lesson.
Reply with JSON only. No text outside the JSON.
RULES
- Exactly 6 or 7 species in total.
- Exactly 3 to 5 trophic stages (level 0 = plant/producer, then up).
- Include "<animal>" itself, at whatever level it really belongs to.
- Include at least one producer at level 0. Include one decomposer.
- IMPORTANT: at least one animal must eat TWO different things. This matters:
  it is what lets a learner test whether a predator survives losing one food.
- Use common English names a 13-year-old knows. No scientific names.
FIELDS
- level: 0 for producers, 1 for plant-eaters, 2, 3, 4 going up.
- abundance: "very common" | "common" | "uncommon" | "rare".
- uncountable: true for mass nouns like grass, false for countable animals.
- links: "eaten" is the food, "eatenBy" is the one doing the eating.
\end{lstlisting}

\textbf{Emits:} structure only --- species, trophic levels, and who eats whom.

\textbf{How it is checked.} The model is never asked for a number that the simulation depends on.
Diet shares are normalised in code to sum to 1 per predator; predation dependency is capped at
0.7 in total per prey so carrying capacity stays positive; baselines come from the abundance
tier and growth rates from trophic level; carrying capacities are then derived so that the
untouched web sits \emph{exactly} at equilibrium. Unknown species ids, self-loops and arrows into
producers are discarded; a consumer left with no food is attached to a plausible species one
level down. Finally \texttt{ensureDecidable()} guarantees at least one predator with two food
sources --- without one, no two agents can disagree about anything a single run can settle.

\textit{Why this split:} asking a small model for coefficients that must satisfy ``shares sum to
1'' and ``dependency below 1'' produces a web that drifts while nothing is happening, which
destroys the premise that every change is attributable to the learner. Asking it whether a frog
eats grasshoppers is a question it answers well.

\textbf{Fallback:} a built-in seven-species web, with the reason shown.

\subsection*{Feature 2 --- Mira (guide) and Aarav (co-learner)}

\textbf{Does:} produces one scaffolding turn and one challenge turn, and classifies what the
learner typed.

\textbf{Receives:} the current web, the learner model, her utterance, and the single proposition
both must address --- computed in code so the two predictions are always comparable.
\textbf{It does not receive the simulation result.}

\begin{lstlisting}
- Mira is a warm guide. She explains and gives confidence. She is the ONLY one
  who may say what to do next. At most 2 short sentences.
- Aarav is a classmate the same age. He is also unsure of things.
  Aarav must DISAGREE with Mira about this exact question. He should argue the
  simple food-chain view: the predator has lost its food, so it will die out.
  He may well be wrong -- that is fine and expected.
CRITICAL RULES
- You do NOT know the simulation result. Do not pretend to.
- Never mention populations as exact numbers. You have not run anything yet.
- Never tell the learner she is wrong. Never use the words test, marks, score.
- Mira and Aarav must address the SAME question.
\end{lstlisting}

\textbf{How it is checked.} Every field is length-clamped and enumerations are constrained to
known values. Template placeholders (``\ldots'', ``\texttt{<text>}'') are rejected as empty, since
a model handed a JSON shape will sometimes return the shape verbatim. Aarav's stance is
\emph{overwritten in code} after the call: it is scheduled deterministically, never left to the
model, so a disagreement can be demonstrated on demand and his wrongness rate stays a
pedagogical parameter rather than a lottery. Because these two agents cannot see the simulator,
there is nothing to fact-check --- their honesty is structural.

\textbf{Fallback:} a recorded exchange that still poses the real disagreement.

\subsection*{Feature 3 --- Dr.\ Sen (evidence checker)}

\textbf{Does:} reports what the numbers show, says who was right, and names the assumption the
result depends on. The only agent that sees the simulation.

\textbf{Receives:} a compact JSON brief of the simulator output, both agents' stated predictions,
and the learner's choice.

\textbf{Emits:} prose plus a per-claim verdict carrying, for each claim, the species it is about
and a direction (up / down / flat).

\textbf{How it is checked --- the strictest check in the system.} Every directional claim is
compared against the simulator: a claim that a species went up when it fell, or that names a
species not in the web, \emph{fails}. The prose is then scanned for numbers, and figures that
correspond to nothing the simulator produced count against it. On failure the model's text is
\textbf{discarded entirely} and replaced by a summary generated directly from the simulation, so
a confident falsehood never reaches the learner. The replacement says which of the two situations
occurred, because ``the helpers are offline'' and ``the helper's reply did not match the result''
are different facts.

\subsection*{Feature 4 --- the doubt chat}

\textbf{Does:} answers free-form questions. Deliberately outside the council pipeline: a learner
who simply wants to know what a decomposer does should not have to commit a prediction and watch
two agents argue to find out.

\textbf{How it is checked.} Mira is again blind to the simulation, so she cannot spoil an
experiment the learner has not run. She is instructed never to state an outcome as fact, and any
web edit she proposes is \emph{confirmed by the learner} before it is applied --- a doubt becomes
an experiment rather than a correction. Proposed edits naming unknown species are dropped.

\textbf{Fallback:} an answer assembled from the web itself (what the named species eats and what
eats it), so it is still true and still about her actual question.

\subsection*{Failure handling across all four features}

\begin{center}
\begin{tabular}{@{}llp{7.4cm}@{}}
\toprule
\textbf{Cause} & \textbf{Learner sees} & \textbf{Rationale}\\
\midrule
All slots 429/503 & ``very busy right now'' & Telling her the helpers are gone when they are
merely loaded sends her away for the wrong reason.\\
Replied, but no JSON & ``got muddled'' & Not an outage; the experiments are unaffected and she
is told so.\\
Unreachable & ``can't be reached'' & The genuine offline case.\\
\bottomrule
\end{tabular}
\end{center}

Three model slots are configured, each with its own model; a slot that returns an unusable reply
fails over exactly as a 503 does, because otherwise one chatty model masks every working slot
behind it. In all cases the simulator keeps running, so the activity never stops.

%====================================================================
\section{D7 \quad Pedagogical rationale}

\subsection*{How the design supports learning}

The activity is built on \emph{commitment before evidence}. The learner ticks who she thinks
changes and commits a prediction before anything runs. This is what converts the simulation from
an animation into a test of her own model: a result is only informative if she had a prior.

Scaffolding is adaptive along three channels, as required for a dynamically adaptive product:
her \emph{responses} (misconception confidence moves on each answer), her \emph{requests for
help}, and her \emph{interaction} (self-initiated web edits lower the support level). These feed
a scaffold level that changes how much Mira explains.

Feedback never corrects. When a misconception fires, she is given something to watch and the
simulation shows her. The cascade narrative is generated in code from the simulator output, so
what she reads can never drift from what the numbers did.

\subsection*{Answer to the P3 design challenge}

\textit{How will the system prevent conflicting agent advice from confusing learners while
keeping learner agency?}

Three invariants, enforced in code:

\begin{enumerate}[leftmargin=*,itemsep=2pt]
  \item \textbf{At most one open disagreement at a time.} While one is unresolved, Aarav's stance
        is forced to ``probe''. Confusion comes from \emph{stacked} unresolved conflicts, not
        from one.
  \item \textbf{Agents may disagree about predictions. They may never disagree about
        instructions.} Only Mira may set the next action. The learner is never told two different
        things to do.
  \item \textbf{Every disagreement must be decidable by one simulator run.} If the conflict
        detector cannot map it to a runnable experiment, it is suppressed and downgraded to a
        question. There are no unresolvable arguments.
\end{enumerate}

Agency is preserved because the learner is the tiebreaker: she commits a choice \emph{before}
seeing the result, ``they're both wrong'' and ``I'm not sure yet'' are first-class options, and
the resolution is an experiment \emph{she} runs, not a verdict an agent hands down. Every
settled finding is recorded with a ``How we know'' line attributed to her action.

\subsection*{The moment where two agents disagree, and how she experiences it}

Aarav says the snakes die out; Mira says they drop but survive. These are both defensible, and
the truth depends on one arrow. The learner sees a single \textbf{Two Views} card --- not two
competing walls of chat --- showing both claims side by side and asking her to decide. The
simulation then gives 18.6: Mira was right. Dr.\ Sen names the arrow and sends her to delete it.
She does; the snakes hit 0; Aarav was right after all.

The design pay-off is the final line: \emph{neither agent was wrong about nature --- they were
picturing different webs}. For a learner whose central fear is being wrong a second time, this
reframes error as having assumed something different, not as being stupid. That is the reason the
disagreement exists at all.

\subsection*{How the learner context shaped the design}

\begin{longtable}{@{}p{5.6cm}p{10.0cm}@{}}
\toprule
\textbf{L5 trait} & \textbf{Design response}\\
\midrule
\endhead
Has already failed this topic & No re-teaching from zero; the first turn is diagnostic and starts
from what survived.\\
High anxiety & No score, no timer, no streak, no red cross. A wrong prediction is called ``a guess
worth making''.\\
Low self-efficacy & Aarav is a \emph{peer}, not a third adult, and is visibly wrong about a third
of the time. She watches someone like her be wrong without consequence before risking it
herself.\\
Expects to be told she is wrong & Nobody corrects her. The simulation shows her, and she runs
it.\\
Will not type if graded & Every free-text box is labelled ``not marked''.\\
Protecting against failing twice & ``What we settled'' accumulates all session and is attributed
to her.\\
Reads correction as proof she is bad at science & The disagreement resolves to ``neither was
wrong --- they pictured different webs''.\\
\bottomrule
\end{longtable}

\subsection*{How C1 shaped the design}

C1 did not merely constrain the interface; it changed the architecture. Because the activity has
to work on a poor connection, the simulation had to be local and deterministic --- and once it
is local and deterministic, it becomes available as \emph{ground truth to check the AI against}.
The constraint produced the verification layer. Two further consequences: the agents' replies are
short because the token budget is small, which also suits a 13-year-old who re-reads when
stressed; and the fallbacks had to be genuinely good, because on this learner's connection they
will be seen.

%====================================================================
\section{D8 \quad Evaluation plan}

\subsection*{Assessment of each outcome}

\begin{longtable}{@{}p{0.4cm}p{3.6cm}p{5.4cm}p{5.2cm}@{}}
\toprule
\textbf{\#} & \textbf{Outcome} & \textbf{Assessment item} & \textbf{Success threshold}\\
\midrule
\endhead
1 & Predict non-adjacent effects & Task 1 multi-select, repeated post-activity on a
    \emph{different} web & $\geq$60\% of learners name at least one non-adjacent species, from a
    $\leq$20\% pre-activity baseline\\
2 & Explain predator survival & Task 3 free text, three-level rubric & $\geq$50\% reach level 3\\
3 & Identify the deciding assumption & Task 4: edit an arrow to flip the outcome &
    $\geq$40\% flip it within three edits\\
\bottomrule
\end{longtable}

\textbf{Rubric for outcome 2.} Level 1: no reference to feeding links. Level 2: names the
alternate prey. Level 3: explains that survival \emph{depends on} the alternate link existing.

\subsection*{Log data collected}

Pre/post ticks with hop distance from the removed species; committed predictions versus outcomes;
free-text explanations with rubric score; Two Views choice distribution (specifically, how often
learners simply side with the adult rather than the peer --- a design risk worth measuring);
number and timing of self-initiated web edits; length of the settled list at exit; abandonment
point; and, for this persona particularly, the \textbf{free-text submission rate}, which tests
whether the ``not marked'' framing actually works.

\subsection*{What would count as success}

Beyond the thresholds above, three behavioural signals: a learner edits the web without being
told to; the settled list is non-empty at exit; and the Two Views choices are not overwhelmingly
``Mira is right'', which would indicate deference to authority rather than reasoning. A measured
gain on outcome 1 with a near-zero free-text rate would count as a partial failure --- it would
mean the content worked but the affective design did not.

%====================================================================
\section*{AI-use log and reflection}
\addcontentsline{toc}{section}{AI-use log and reflection}

\subsection*{Log}

\textit{Insert the share links or exports of every AI chat used, in the order used.}

\begin{enumerate}[leftmargin=*,itemsep=1pt]
  \item Topic and card selection; comparison of P3-T4 against P1-T3 on available free tooling.
  \item Product specification (\texttt{PRD.md}): agent roles, collaboration protocol, conflict
        invariants.
  \item Implementation: simulator, agent pipeline, web generation, authentication, interface.
  \item Debugging the model layer; diagnosis of the fallback behaviour.
  \item This report.
\end{enumerate}

\subsection*{Reflection}

AI was used throughout, and the most useful thing I learned was \emph{where not to use it}.

The first version asked the model for the simulation coefficients. The result looked plausible
and was pedagogically useless: the web drifted while nothing was happening, so a learner could
not attribute any change to her own action. Moving every number into code and leaving the model
only the structure --- which species, who eats whom --- fixed it. The same division of labour
runs through the whole system: the model proposes, the code disposes.

The second lesson was that an agent's \emph{ignorance} can be a design asset. Mira and Aarav are
deliberately not shown the simulation result, which is what makes their disagreement honest
rather than theatre, and which means there is nothing to fact-check. Only Dr.\ Sen sees the
numbers, and he is therefore the only one who needs verification --- which is implemented, and
which discards his text outright when it contradicts the simulator.

The third lesson came from a bug. The application kept reporting that the agents were offline
when they were not: one model was replying with paragraphs of thinking-aloud instead of JSON, and
my code treated ``answered unusably'' as success, so a working provider behind it was never
tried. Two things were wrong --- the failover logic, and the message, which told the learner
something untrue. Both were fixed, and the lesson generalises: when an AI feature degrades, the
system must be specific about \emph{how}, because the honest message and the convenient message
are rarely the same.

What I would not delegate again is the feedback text for the misconceptions. The generated
versions were fluent and subtly corrective in tone --- exactly wrong for a learner who stops when
she feels judged. Those strings are written by hand, in code, and they do not change between runs.

\end{document}
